\documentclass[11pt,a4paper]{article}
\usepackage[margin=1in]{geometry}
\usepackage{times}
\usepackage{latexsym}
\usepackage[colorlinks=true,citecolor=blue,urlcolor=blue]{hyperref}
\usepackage{natbib}
\renewcommand{\citename}[1]{#1}
\usepackage{graphicx}
\usepackage{booktabs}
\usepackage{amsmath}
\usepackage{multirow}
\usepackage{xcolor}

\title{The Training Lottery: Why Error Injection Destabilizes\\Embodied Planning Agents Across Architectures}

\author{Anonymous}

\begin{document}
\maketitle

\begin{abstract}
Error injection is a widely used data augmentation strategy for training robust LLM agents.
Prior work on tool-use agents (PALADIN) reports large gains from taxonomy-guided error injection paired with recovery demonstrations.
We test whether this approach transfers to embodied planning by designing a 4-operator error taxonomy for ALFWorld and training Qwen3-8B and LLaMA-3.1-8B under three conditions: Vanilla (clean data), TaxEI (taxonomy-labeled injection), and Random (unlabeled injection).
Across 5 seeds on Qwen3-8B, both injection methods degrade performance identically ($58.3\%$ vs.\ Vanilla $67.3\%$, $-9.0$pp, Cohen's $d{=}1.58$), and taxonomy labels provide zero benefit (TaxEI $=$ Random $= 58.3\%$).
Critically, within-condition seed variance (19--23pp range) dominates between-condition differences at every injection dose, and 1-in-5 seeds trigger catastrophic failure (${\sim}47\%$) absent from clean training.
Cross-model replication on LLaMA-3.1-8B confirms the instability pattern: TaxEI SD${=}8.6$pp vs.\ Random SD${=}2.1$pp, with 1/3 seeds showing catastrophic degradation.
Recovery evaluation reveals that models exhibit \emph{resilience} (passively ignoring errors) rather than \emph{recovery} (detecting and correcting errors), with 0\% detection rate across all conditions.
Scaling training data from 501 to 2500 trajectories narrows but does not eliminate the harm ($-5.7$pp, NS).
Per-operator effects observed at single-seed resolution vanish under multi-seed validation.
These findings establish that error injection outcomes in embodied planning are dominated by the \textbf{training lottery}---stochastic seed-dependent variation---rather than by injection method, dose, or taxonomy design, posing a methodological warning for the field.
\end{abstract}

%-------------------------------------------------------------------
\section{Introduction}
\label{sec:intro}

Error injection---deliberately introducing failure modes into training data so models learn to handle them---is an intuitive augmentation strategy for LLM agents.
PALADIN \cite{paladin2025} demonstrated its potential on tool-use agents, achieving a +57pp Recovery Rate gain on ToolBench through taxonomy-guided error injection combined with inference-time retrieval augmentation.
A natural question follows: does this approach transfer to other agent domains?

We investigate this question for embodied planning, where agents execute multi-step physical manipulation tasks.
We design TaxEI, a 4-operator error taxonomy for ALFWorld \cite{alfworld2021}, and conduct a comprehensive evaluation across two model architectures (Qwen3-8B, LLaMA-3.1-8B-Instruct), five random seeds, five injection doses, and multiple evaluation protocols.

Our central finding is \emph{negative but methodologically informative}: error injection consistently harms embodied planning agents, but the dominant source of performance variation is not the injection method, dose, or taxonomy design---it is the \textbf{training lottery}, i.e., stochastic seed-dependent outcomes.
Within-condition seed variance (19--23pp) exceeds between-condition mean differences (9pp) at every tested dose.
Both structured (TaxEI) and unstructured (Random) injection produce identical mean degradation ($-9.0$pp) on Qwen3-8B, proving taxonomy labels provide zero benefit.
Cross-model replication on LLaMA-3.1-8B shows the instability is architecture-general for structured injection.

This work contributes: (1) a \emph{training lottery} methodological warning---single-seed evaluations of error injection are fundamentally unreliable; (2) the first cross-domain transfer test of taxonomy-guided error injection, revealing task-structural boundary conditions; (3) evidence that observed error ``recovery'' in embodied agents is passive resilience, not active error correction; and (4) a complete dose-response profile showing no safe injection threshold exists.

%-------------------------------------------------------------------
\section{Related Work}
\label{sec:related}

\paragraph{Error Injection for LLM Agents.}
PALADIN \cite{paladin2025} uses ToolScan to generate taxonomy-guided error-recovery pairs for tool-use agents, achieving large gains on ToolBench.
However, PALADIN's design entangles SFT error injection with inference-time retrieval of 55+ exemplars; removing retrieval reduces gains by ${\sim}30$pp.
EEF \cite{eef2025} takes the inverse approach on WebShop, extracting positive signal from failure trajectories rather than injecting errors, suggesting the injection direction itself may be suboptimal for planning.

\paragraph{Training Instability and Seed Sensitivity.}
McCoy et al.\ \citeyearpar{mccoy2020} and Mosbach et al.\ \citeyearpar{mosbach2021} document that fine-tuning outcomes for pre-trained models are highly seed-dependent, especially on small datasets.
D'Amour et al.\ \citeyearpar{damour2022} formalize this as \emph{underspecification}: models that perform identically on average can diverge dramatically on specific inputs depending on training randomness.
Our work extends these findings to error-injection augmentation, showing that the training lottery effect is amplified by injected perturbations.

\paragraph{Embodied Planning Agents.}
ALFWorld \cite{alfworld2021} is a standard benchmark for text-based embodied planning.
ReAct \cite{yao2023react} and Reflexion \cite{shinn2023reflexion} improve agent performance through inference-time reasoning, not training-time augmentation.
Our work complements these by showing that training-time error injection does not confer the same benefits.

\paragraph{Adversarial Training and Curriculum Learning.}
Adversarial training \cite{madry2018} and curriculum learning \cite{bengio2009} both modify training distributions to improve robustness.
Our error injection can be viewed as a form of curriculum learning that fails when the curriculum's error signals are diffuse and unobservable.

%-------------------------------------------------------------------
\section{Method}
\label{sec:method}

\subsection{Error Taxonomy for Embodied Planning}

We design a 4-operator error taxonomy targeting distinct failure modes in ALFWorld:

\begin{itemize}
\item \textbf{P1} (\texttt{step\_order\_swap}): Permute the ordering of action steps.
\item \textbf{P2} (\texttt{sub\_goal\_deletion}): Remove sub-goals from the plan.
\item \textbf{C1} (\texttt{context\_truncation}): Truncate environmental context information.
\item \textbf{C4} (\texttt{observation\_masking}): Mask observation feedback from the environment.
\end{itemize}

These operators span a spectrum from local/observable (C4) to global/cascading (C1, P2) perturbations.
All operators pass injection legality validation (100\% structurally valid) and recovery quality checks (97.5\%).

\subsection{Training Setup}

We generate 501 successful trajectories on ALFWorld using GPT-5.2, producing 4,671 injected variants across the 4 operators.
Three training conditions are evaluated:

\begin{itemize}
\item \textbf{Vanilla}: Clean trajectories only (4,358 samples).
\item \textbf{TaxEI}: Clean + taxonomy-labeled injections (9,029 samples).
\item \textbf{Random}: Clean + unlabeled injections (9,029 samples, identical data, generic error tags instead of taxonomy labels).
\end{itemize}

All conditions are trained via SFT on Qwen3-8B with 5 seeds (42, 0, 1, 7, 123).
Cross-model validation uses LLaMA-3.1-8B-Instruct with 3 seeds (42, 0, 123).

\subsection{Evaluation Protocols}

\paragraph{Clean Evaluation.} Standard ALFWorld task completion on 47 test tasks.
\paragraph{Recovery Evaluation.} Tasks with round-robin error injection (P1/P2/C1/C4) at evaluation time, measuring whether training-time injection improves robustness to eval-time errors.
\paragraph{Detection Rate.} Whether models explicitly flag injected errors in their output.

%-------------------------------------------------------------------
\section{Experiments}
\label{sec:experiments}

\subsection{Main Results: The Training Lottery Dominates}
\label{sec:main}

Table~\ref{tab:main_5seed} presents the core 5-seed results on Qwen3-8B.
Both injection methods produce \emph{identical} mean degradation: TaxEI $= $ Random $= 58.3\%$ vs.\ Vanilla $67.3\%$ ($-9.0$pp).
Taxonomy labels provide zero benefit (Welch $t{=}0.000$, $p{=}1.000$).

The defining feature is seed variance: within each condition, the range spans 19--28pp (e.g., TaxEI: 46.8\%--61.7\%).
Both injection methods produce exactly one catastrophic outlier (${\sim}47\%$), absent from Vanilla (minimum 61.7\%).
Welch $t$ for Vanilla vs.\ TaxEI yields $p{=}0.039$ uncorrected, but Bonferroni correction across 3 comparisons ($\alpha{=}0.017$) renders all comparisons non-significant.
Cohen's $d{=}1.58$ indicates a large effect size, but the study is underpowered at $n{=}5$ (power $= 45.2\%$; $n{=}9$ required for 80\% power).

\begin{table}[t]
\centering
\small
\caption{5-seed clean success rate on Qwen3-8B (47 ALFWorld tasks). TaxEI and Random produce identical mean degradation. Within-condition ranges exceed between-condition differences.}
\label{tab:main_5seed}
\begin{tabular}{lcccccc}
\toprule
\textbf{Cond.} & \textbf{s42} & \textbf{s0} & \textbf{s1} & \textbf{s7} & \textbf{s123} & \textbf{Mean$\pm$SD} \\
\midrule
Vanilla & 61.7 & 68.1 & 66.0 & 74.5 & 61.7 & 67.3$\pm$4.7 \\
TaxEI & 59.6 & 46.8 & 61.7 & 61.7 & 55.3 & 57.0$\pm$6.2 \\
Random & 63.8 & 55.3 & 51.1 & 55.3 & 66.0 & 58.3$\pm$6.4 \\
\bottomrule
\end{tabular}
\end{table}

\subsection{Recovery Evaluation: Resilience, Not Recovery}
\label{sec:recovery}

Table~\ref{tab:recovery_5seed} presents 5-seed recovery evaluation results under eval-time error injection.
Vanilla achieves $66.4\%{\pm}5.1$pp, TaxEI $57.0\%{\pm}6.2$pp, Random $58.3\%{\pm}6.4$pp.
Recovery rates closely match clean evaluation rates (delta $< 2$pp across all conditions), confirming that models \emph{passively ignore} eval-time perturbations rather than actively detecting and recovering from them.
Detection rate is 0\% for all conditions on Qwen3-8B.

\begin{table}[t]
\centering
\small
\caption{5-seed recovery rate under injected evaluation (47 tasks). Recovery $\approx$ clean rates across all conditions, indicating resilience rather than recovery.}
\label{tab:recovery_5seed}
\begin{tabular}{lcccccc}
\toprule
\textbf{Cond.} & \textbf{s42} & \textbf{s0} & \textbf{s1} & \textbf{s7} & \textbf{s123} & \textbf{Mean$\pm$SD} \\
\midrule
Vanilla & 61.7 & 68.1 & 66.0 & 74.5 & 61.7 & 66.4$\pm$5.1 \\
TaxEI & 59.6 & 46.8 & 61.7 & 61.7 & 55.3 & 57.0$\pm$6.2 \\
Random & 63.8 & 55.3 & 51.1 & 55.3 & 66.0 & 58.3$\pm$6.4 \\
\bottomrule
\end{tabular}
\end{table}

\subsection{Recovery Prompting: Marginal Effect}
\label{sec:recovery_prompt}

We test whether few-shot recovery prompting at inference time can compensate for the lack of training-time error awareness.
Table~\ref{tab:recovery_prompt} shows results with 3-shot recovery exemplars (seed=42).
Prompting provides only marginal improvement (+2.1pp for Vanilla and TaxEI, 0.0pp for Random) and raises detection rate from 0\% to only 8.5\%.
Even when given explicit recovery knowledge, models predominantly rely on resilience rather than active error correction.

\begin{table}[t]
\centering
\small
\caption{Recovery prompting results (seed=42, 47 tasks). Few-shot exemplars provide marginal improvement.}
\label{tab:recovery_prompt}
\begin{tabular}{lcccc}
\toprule
\textbf{Condition} & \textbf{Base} & \textbf{3-shot} & \textbf{$\Delta$pp} & \textbf{Det.\ rate} \\
\midrule
Vanilla & 61.7\% & 63.8\% & +2.1 & 8.5\% \\
TaxEI & 59.6\% & 61.7\% & +2.1 & 8.5\% \\
Random & 63.8\% & 63.8\% & 0.0 & 6.4\% \\
\bottomrule
\end{tabular}
\end{table}

\subsection{Cross-Model Validation: LLaMA-3.1-8B}
\label{sec:crossmodel}

To test whether the training lottery is architecture-specific, we replicate the core experiment on LLaMA-3.1-8B-Instruct with 3 seeds (Table~\ref{tab:llama}).
TaxEI shows high variance (SD${=}8.6$pp) with 1/3 seeds triggering catastrophic degradation (s123: $59.6\%$ vs.\ Vanilla $70.2\%$, $\Delta{=}-10.6$pp).
Random injection remains stable (SD${=}2.1$pp) and mildly beneficial (mean $76.6\%$ vs.\ Vanilla $73.1\%$).

Table~\ref{tab:crossmodel_summary} summarizes cross-model consistency.
The direction of TaxEI harm holds across architectures, but key patterns \emph{diverge}: Random injection is harmful on Qwen but beneficial on LLaMA; TaxEI variance is lower than Random on Qwen but $4.1\times$ higher on LLaMA; detection rate is 0\% on Qwen but 25.5\% on LLaMA.
The only consistent finding is that structured injection (TaxEI) produces at least one catastrophic seed on both models.

\begin{table}[t]
\centering
\small
\caption{LLaMA-3.1-8B recovery rate (injected eval, 47 tasks, 3-seed). TaxEI shows high variance with catastrophic seed s123.}
\label{tab:llama}
\begin{tabular}{lcccc}
\toprule
\textbf{Cond.} & \textbf{s42} & \textbf{s0} & \textbf{s123} & \textbf{Mean$\pm$SD} \\
\midrule
Vanilla & 74.5 & 74.5 & 70.2 & 73.1$\pm$2.5 \\
TaxEI & 74.5 & 74.5 & 59.6 & 69.5$\pm$8.6 \\
Random & 76.6 & 74.5 & 78.7 & 76.6$\pm$2.1 \\
\bottomrule
\end{tabular}
\end{table}

\begin{table}[t]
\centering
\small
\caption{Cross-model comparison. Only the direction of TaxEI harm is consistent; magnitude, variance patterns, and detection rates diverge.}
\label{tab:crossmodel_summary}
\begin{tabular}{lccc}
\toprule
\textbf{Metric} & \textbf{Qwen3-8B} & \textbf{LLaMA-3.1-8B} & \textbf{Consistent?} \\
\midrule
TaxEI vs Van.\ $\Delta$ & $-9.0$pp & $-3.6$pp & Direction only \\
Rand.\ vs Van.\ $\Delta$ & $-9.0$pp & $+3.5$pp & \textbf{No} \\
TaxEI SD / Rand.\ SD & $0.75\times$ & $4.1\times$ & \textbf{No} \\
Catastrophic seeds & 1/5 & 1/3 & Yes \\
Detection rate & 0\% & 25.5\% & \textbf{No} \\
\bottomrule
\end{tabular}
\end{table}

\subsection{Scale Ablation: 501 vs.\ 2500 Trajectories}
\label{sec:scale}

We test whether insufficient data scale explains the negative result by expanding training data from 501 to 2500 trajectories (Table~\ref{tab:scale}).
At $5\times$ scale, the TaxEI--Vanilla gap narrows from $-9.0$pp to $-5.7$pp but remains negative and non-significant (Welch $t$, $p{=}0.18$).
This experiment is inherently limited: ALFWorld has only 134 unique tasks, so 2500 trajectories represent ${\sim}18.7$ repetitions per task---scaling repetition, not task diversity.
By contrast, PALADIN's 50K+ trajectories span 5000+ distinct tool-use scenarios (Table~\ref{tab:scale_feasibility}).
Meaningful scale-matched comparison is infeasible within ALFWorld's task space.

\begin{table}[t]
\centering
\small
\caption{Scale ablation (3-seed). Gap narrows but persists; Bonferroni NS at both scales.}
\label{tab:scale}
\begin{tabular}{lcccc}
\toprule
\textbf{Scale} & \textbf{Vanilla} & \textbf{TaxEI} & \textbf{$\Delta$pp} & \textbf{Welch $p$} \\
\midrule
501 traj & 67.3$\pm$4.7 & 58.3$\pm$6.5 & $-9.0$ & 0.039* \\
2500 traj & 70.2$\pm$3.8 & 64.5$\pm$5.9 & $-5.7$ & 0.18 \\
\bottomrule
\multicolumn{5}{l}{\footnotesize * Uncorrected; Bonferroni NS at both scales.}
\end{tabular}
\end{table}

\begin{table}[t]
\centering
\small
\caption{Scale feasibility: ALFWorld's limited task diversity makes PALADIN-scale comparison invalid.}
\label{tab:scale_feasibility}
\begin{tabular}{lccc}
\toprule
\textbf{Dataset} & \textbf{Unique tasks} & \textbf{50K $=$ reps/task} & \textbf{Diversity} \\
\midrule
PALADIN ToolBench & 5000+ & ${\sim}10\times$ & High \\
ALFWorld (ours) & 134 & ${\sim}373\times$ & Low \\
\bottomrule
\end{tabular}
\end{table}

\subsection{Per-Operator Analysis: Single-Seed Artifacts}
\label{sec:peroperator}

A single-seed (s42) analysis suggested a local/global split: C4 (observation masking, local) showed TaxEI advantage (+9.1pp) while C1/P2 (global/cascading) showed disadvantage ($-8.3$pp each).
We validate this with 5-seed pooled data ($N{=}55$--$60$ per operator, Table~\ref{tab:peroperator}).
The C4 advantage collapses from $+9.1$pp to $+1.8$pp (Fisher $p{=}0.52$).
No operator reaches significance (all $p > 0.35$).
Per-seed operator rankings are unstable---some seeds show C4 as worst, others as best.
The apparent local/global pattern was an artifact of $N{=}11$--$12$ sampling noise at single-seed resolution.

\begin{table}[t]
\centering
\small
\caption{Per-operator TaxEI--Vanilla recovery gap (5-seed pooled). All effects are non-significant; the single-seed local/global pattern does not replicate.}
\label{tab:peroperator}
\begin{tabular}{llcccc}
\toprule
\textbf{Operator} & \textbf{Type} & \textbf{TaxEI} & \textbf{Vanilla} & \textbf{$\Delta$pp} & \textbf{Fisher $p$} \\
\midrule
C4 obs.\ mask & Local & 54.5 & 52.7 & +1.8 & 0.52 \\
P1 step swap & Structural & 56.7 & 57.5 & $-0.8$ & 0.89 \\
C1 ctx.\ trunc. & Global & 58.3 & 61.5 & $-3.2$ & 0.38 \\
P2 goal del. & Global & 51.7 & 54.4 & $-2.7$ & 0.44 \\
\bottomrule
\end{tabular}
\end{table}

This non-replication is itself evidence for the training lottery thesis: even per-operator effect directions are dominated by seed-dependent stochasticity, not systematic properties of error types.

%-------------------------------------------------------------------
\section{Discussion}
\label{sec:discussion}

\subsection{The Training Lottery as Methodological Warning}

Our results reveal that error injection amplifies the training lottery effect documented by McCoy et al.\ \citeyearpar{mccoy2020} and D'Amour et al.\ \citeyearpar{damour2022}.
Under clean training, Qwen3-8B achieves stable performance (SD${=}4.7$pp, range 61.7--74.5\%).
Both injection methods inflate variance (TaxEI SD${=}6.5$pp, Random SD${=}8.7$pp) and introduce a 1-in-5 risk of catastrophic failure (${\sim}47\%$) absent from Vanilla.
At every tested dose (10\%, 25\%, 50\%, Full), within-condition ranges span 19--23pp.

This has direct implications for how error injection methods are evaluated.
PALADIN's reported gains use a single training configuration; our results suggest that such reports may be sensitive to seed selection.
We recommend that future error injection studies report $\geq$5 seed results with explicit variance analysis.

\subsection{Why Error Injection Fails in Embodied Planning}

We hypothesize that the divergence from PALADIN's positive results stems from structural differences between tool-use and embodied planning:

\begin{itemize}
\item \textbf{Error observability}: Tool calls return explicit error codes; embodied actions produce syntactically valid but semantically wrong observations.
\item \textbf{Recovery locality}: Failed API calls can be locally retried; a wrong embodied action invalidates the entire downstream plan.
\item \textbf{Feedback granularity}: Tool-use provides per-call success/fail signals; ALFWorld provides only sparse task-completion reward.
\end{itemize}

These structural differences mean that injected errors in embodied planning are indistinguishable from noise, corrupting the action distribution rather than providing useful curriculum signal.

\subsection{Resilience vs.\ Recovery}

Recovery evaluation rates match clean rates within 2pp across all conditions (Section~\ref{sec:recovery}), with 0\% detection rate on Qwen3-8B.
Models succeed at the same base rate regardless of eval-time error injection---they \emph{passively ignore} perturbations.
This is resilience (noise tolerance from task competence), not recovery (error detection + correction).

Recovery prompting (Section~\ref{sec:recovery_prompt}) provides only +2.1pp improvement and raises detection to 8.5\%, suggesting that error awareness is difficult to elicit in embodied planning even with explicit prompting.
This aligns with the structural hypothesis: without discrete error signals in the observation space, there is nothing for the model to detect.

LLaMA-3.1-8B shows 25.5\% detection rate, in contrast to Qwen's 0\%.
This model-dependent divergence further complicates any universal claim about error awareness in embodied planning and suggests that detection capacity may depend on pre-training data composition rather than fine-tuning augmentation.

\subsection{Limitations}

\paragraph{Statistical power.} At $n{=}5$ seeds, power is 45.2\% for the observed effect at Bonferroni-corrected $\alpha{=}0.017$. We report $n{=}9$ as the threshold for 80\% power. All per-comparison $p$-values are Bonferroni NS.

\paragraph{Cross-model scope.} LLaMA replication uses only 3 seeds. The inconsistent Random injection effect (harmful on Qwen, beneficial on LLaMA) limits cross-model generalization.

\paragraph{Task diversity.} ALFWorld has 134 unique tasks and 47 test tasks---a constraint shared with all prior work on this benchmark but limiting absolute statistical resolution. Scale ablation is bounded by task diversity, not trajectory count.

\paragraph{Data volume confound.} Vanilla trains on 4,358 samples vs.\ 9,029 for TaxEI/Random. Dose-response ablation at 10\% injection (fixed 4,358 total) confirms harm from injection content, not volume, but the confound remains in the full-scale comparison.

\paragraph{PALADIN comparison caveats.} The comparison involves multiple confounds beyond task structure: base model, training scale (9K vs.\ 50K+), and inference-time retrieval augmentation (${\sim}30$pp contribution to PALADIN's gains). We test error injection in its most isolated form (pure SFT, no retrieval).

%-------------------------------------------------------------------
\section{Conclusion}
\label{sec:conclusion}

We present the first cross-domain transfer test of taxonomy-guided error injection for LLM agents, from tool-use to embodied planning.
The transfer fails: both structured and unstructured injection identically degrade Qwen3-8B performance by 9.0pp, and taxonomy labels provide zero benefit.
Cross-model validation on LLaMA-3.1-8B confirms that structured injection introduces training instability across architectures.

The dominant finding is methodological: error injection amplifies the \emph{training lottery}, producing within-condition variance (19--23pp) that exceeds between-condition mean differences (9pp) at every dose.
Single-seed evaluations of error injection methods are fundamentally unreliable.
Per-operator effects that appear systematic at $N{=}11$ dissolve into noise at $N{=}55$.
Recovery evaluation reveals resilience, not recovery, with 0\% error detection.

These results establish task-structural boundary conditions for error injection: the approach requires discrete, observable, locally recoverable errors---properties present in tool-use but absent in embodied planning.
We urge the community to adopt multi-seed evaluation protocols for error injection studies and to distinguish resilience from recovery when evaluating agent robustness.

\begin{thebibliography}{20}

\bibitem[Shridhar et~al.(2021)]{alfworld2021}
Mohit Shridhar, Xingdi Yuan, Marc-Alexandre C{\^o}t{\'e}, Yonatan Bisk, Adam Trischler, and Matthew Hausknecht. 2021.
\newblock ALFWorld: Aligning Text and Embodied Environments for Interactive Learning.
\newblock In \emph{ICLR}.

\bibitem[{PALADIN Authors}(2025)]{paladin2025}
{PALADIN Authors}. 2025.
\newblock PALADIN: Taxonomy-Guided Error Injection for Tool-Use Agent Training.
\newblock \emph{arXiv:2509.25238}.

\bibitem[{EEF Authors}(2025)]{eef2025}
{EEF Authors}. 2025.
\newblock Extracting Experience from Failures for LLM Agents.
\newblock \emph{arXiv:2504.13145}.

\bibitem[McCoy et~al.(2020)]{mccoy2020}
R.~Thomas McCoy, Junghyun Min, and Tal Linzen. 2020.
\newblock BERTs of a feather do not generalize together: Large variability in generalization across models with similar test set performance.
\newblock In \emph{Proceedings of the Third BlackboxNLP Workshop}.

\bibitem[Mosbach et~al.(2021)]{mosbach2021}
Marius Mosbach, Maksym Andriushchenko, and Dietrich Klakow. 2021.
\newblock On the Stability of Fine-tuning BERT: Misconceptions, Explanations, and Strong Baselines.
\newblock In \emph{ICLR}.

\bibitem[D'Amour et~al.(2022)]{damour2022}
Alexander D'Amour, Katherine Heller, Dan Moldovan, et~al. 2022.
\newblock Underspecification Presents Challenges for Credibility in Modern Machine Learning.
\newblock \emph{JMLR}, 23(226):1--61.

\bibitem[Yao et~al.(2023)]{yao2023react}
Shunyu Yao, Jeffrey Zhao, Dian Yu, Nan Du, Izhak Shafran, Karthik Narasimhan, and Yuan Cao. 2023.
\newblock ReAct: Synergizing Reasoning and Acting in Language Models.
\newblock In \emph{ICLR}.

\bibitem[Shinn et~al.(2023)]{shinn2023reflexion}
Noah Shinn, Federico Cassano, Ashwin Gopinath, Karthik Narasimhan, and Shunyu Yao. 2023.
\newblock Reflexion: Language Agents with Verbal Reinforcement Learning.
\newblock In \emph{NeurIPS}.

\bibitem[Madry et~al.(2018)]{madry2018}
Aleksander Madry, Aleksandar Makelov, Ludwig Schmidt, Dimitris Tsipras, and Adrian Vladu. 2018.
\newblock Towards Deep Learning Models Resistant to Adversarial Examples.
\newblock In \emph{ICLR}.

\bibitem[Bengio et~al.(2009)]{bengio2009}
Yoshua Bengio, J{\'e}r{\^o}me Louradour, Ronan Collobert, and Jason Weston. 2009.
\newblock Curriculum Learning.
\newblock In \emph{ICML}.

\end{thebibliography}

\end{document}
